\chap[passkeys-theory] Passkeys

\nl The Internet has become an integral part of daily life for most of the world’s population. In December 2025, approximately {\sbf 6.047 billion people} (73.2\% of the global population) were {\sbf Internet users}, representing a 14\% increase compared to 2020~\cite[internet-usage-stats]. This growth, along with an increasing reliance on online services, has been accompanied by an increase in cyberattacks and personal data breaches.

\nl This development highlights a fundamental weakness of password-based authentication. According to the Verizon Data Breach Investigations Report, approximately {\sbf 88\% of breaches} in 2025 involved {\sbf stolen credentials}~\cite[verizon-report-2025]. Passwords are therefore inherently vulnerable to theft, reuse, and brute-force attacks.

\nl In response, organizations have {\sbf increasingly adopted multi-factor authentication (MFA}), which supplements passwords with additional verification factors such as one-time codes, authentication applications, or trusted devices. The adoption of MFA has grown significantly; for example, {\sbf 83\% of organizations required MFA} for all access to IT in 2024~\cite[jumpcloud-sme-2024]. In the Czech Republic, the use of MFA by organizations is {\sbf required by law} since 2025. Despite variations across surveys, it is clearly observable that {\sbf MFA has become a widely used safety measure}. 

\sec[passkeys] The Advantages of Passkeys

Passkeys -- standardized by the FIDO Alliance -- have emerged as a response to the limitations of traditional multi-factor authentication (MFA). Many conventional MFA methods are considered to be insufficiently secure, while simultaneously introducing unnecessary complexity to the authentication process, thus worsening the user experience.

\nl The FIDO Alliance {\sbf defines a passkey} as
 ``\dots authentication {\sbf credential} based on FIDO standards that allows a user to {\sbf log into apps and websites} with the same process they use to unlock their device (biometrics, PIN or pattern). Passkeys rely on {\sbf asymmetric cryptography}, where the private key is securely stored on the user’s authenticator, while the corresponding public key is registered with the service ''.
\urlnote{https://fidoalliance.org/passkeys/}

\nl According to data published by the FIDO Alliance, passkey support has already been {\sbf adopted} by approximately a {\sbf fifth} of the world’s {\sbf top 100 websites} and services~\cite[fido-password-passkey-trends]. Furthermore, major technology providers, including Apple, Google, and Microsoft, have integrated passkey support into their platforms, allowing further support for this authentication method.

\secc[hardware] Hardware security keys

Although passkeys are often implemented as platform authenticators integrated into user devices, the {\sbf FIDO2} framework also supports {\sbf external authenticators} in the form of {\sbf hardware security keys}. These devices {\sbf store cryptographic credentials} on dedicated hardware and {\sbf perform authentication independently} of the host system and only {\sbf upon user interaction}, providing an additional layer of security.

\nl  Currently, the hardware security key market is dominated by several major vendors, most notably {\sbf Yubico} with its Security Key Series. Other notable solutions include the Google Titan Security Key developed by Google and the SecurID authenticators made by RSA Security. In addition, several smaller or less widely cited manufacturers provide compatible FIDO2 security keys, including Feitian Technologies, Swissbit, and TrustKey.

\nl Quantitative estimates of market share are difficult to obtain, as publicly available data are limited and often based on proprietary methodologies. {\sbf Estimates  suggest}, as of April 2026~\cite[yubico-market-share] that {\sbf Yubico accounts for approximately 32\% and RSA Security for 24\%} of the 2FA {\sbf solutions}. However, these figures reflect only the customer base observed by the source platform and should be interpreted as indicative rather than representative of the global market share of hardware security keys.

\nl A distinguishing {\sbf characteristic} of most commercially available hardware security keys is their {\sbf closed-source firmware}. Open-source alternatives do exist but remain rare: the first publicly known example was SoloKeys (developed 2018–2024, now discontinued), followed by OpenSK by Google and Nitrokey. To the best of our knowledge, no open-source C-based implementation of CTAP2 over NFC for a hardware security key was publicly available at the time of writing.

\sec[fido2-framework] The FIDO2 Framework

\nl The {\sbf FIDO2} framework, developed by the FIDO Alliance in collaboration with the W3C, is the current {\sbf industry standard} for {\sbf strong, phishing-resistant and passwordless authentication}. It is an open framework that uses public-key cryptography.

\nl FIDO2 is composed of two complementary core specifications:

\begitems

* {\sbf WebAuthn (Web Authentication):} A W3C standard API that enables web applications to create and use public-key credentials. It defines how a browser or platform (the ``Client'') communicates with a web service (the ``Relying Party'') to register and authenticate users.

* {\sbf CTAP (Client-to-Authenticator Protocol):} A FIDO Alliance specification that defines how a Client communicates with an authenticator (such as a hardware security key or a smartphone) to fulfill WebAuthn requests.

\enditems

The CTAP standard defines two primary versions:

\begitems

* {\sbf CTAP1 (formerly U2F):} Supports legacy FIDO U2F credentials.

* {\sbf CTAP2:} The current standard for FIDO2, which improves user verification methods, most notably by implementing a PIN/User Verification (UV) authentication model.

\enditems
A {\sbf detailed breakdown} of this framework, including a detailed description of the CTAP protocol, can be {\sbf found} in the {\sbf work of Endler}~\cite[endler-thesis], along with a more comprehensive insight into the topic of hardware security keys, including the cryptographic mechanisms that are commonly used. 

\secc[ctap-safety] Safety of CTAP over NFC

Since the protocol relies on {\sbf asymmetric cryptography} and challenge-response authentication, secret credentials are not transmitted directly over the communication channel, significantly reducing the risk of credential theft through passive eavesdropping. However, several attacks targeting FIDO/CTAP authenticators have been described. 

\nl Regarding NFC specifically, there is the possibility of {\sbf Client Impersonation (CI)} and {\sbf API Confusion (AC)} attacks. These methods and their results are described in detail by Casagrande and Antonioli~\cite[casagrande:hal-05406456]. The authors also suggest countermeasures to these attacks, among others, the most notable are different PIN for destructive commands, user interaction for user presence verification (i.e., pressing a button instead of only presenting the authenticator), and improved authenticator visual feedback. 

\nl Although the proposed countermeasures could be integrated without significant implementation complexity, the implementation presented in this thesis follows the CTAP specification and therefore does not introduce additional proprietary protection mechanisms beyond those required by the standard. It is worth noting that despite these exploits being published in 2023, none of the aforementioned methods has been introduced to the CTAP standard.

\secc[ctap-relevancy] Post-Quantum Security of the CTAP Protocol

As quantum computing technology advances, it poses a significant {\sbf long-term threat} to the currently used public-key cryptographic schemes~\cite[enwiki:1349019275]. Although no quantum computer is currently capable of breaking them, the issue is important to acknowledge. Therefore, we provide a {\sbf brief examination} of the {\sbf post-quantum (PQ) security} of the CTAP protocol.

\nl The version of CTAP implemented in this project (CTAP 2.2) {\sbf does not provide native PQ security}, as it relies on conventional public-key cryptographic primitives that could be compromised by sufficiently capable quantum adversaries. Nevertheless, the {\sbf current protocol architecture} is considered to be {\sbf PQ ready}. In particular, its PIN/UV authentication model requires active user verification and physical presence, while the protocol itself can be adapted to stronger cryptographic primitives without a fundamental redesign. Therefore, a full quantum-resistant operation {\sbf could be achieved} by replacing the current asymmetric mechanisms with post-quantum key encapsulation and signature mechanisms, provided that such changes are introduced consistently across both CTAP and WebAuthn within the FIDO2 framework~\cite[cryptoeprint:2022/1029]. The topic of making the CTAP protocol quantum-safe is discussed in more detail in work of {\sbf Kniep}~\cite[kniep-thesis].
